\section{Introduction}
\label{sec:introduction}

Pretrained speech encoders provide representations that can be reused across downstream applications. Self-supervised models such as wav2vec~2.0 learn from unlabeled speech, while models such as Whisper learn through large-scale transcript supervision \citep{baevski2020wav2vec,radford2023whisper}. These approaches offer multiple candidate encoders for languages with limited task-specific supervision. Choosing among them, however, requires more than identifying a strong speech-recognition model: an embedding used to retrieve spoken content must preserve different distinctions from one used to identify speakers. For Sinhala, we investigate which existing representations support these capabilities, how their extraction affects performance, and whether the resulting recommendations generalize across corpora.

Existing benchmarks provide an important foundation for this investigation. ML-SUPERB evaluates multilingual speech representations through recognition and language-identification tasks, and its successor studies the effects of downstream configurations \citep{shi2023mlsuperb,shi2024mlsuperb2}. The Massive Audio Embedding Benchmark (MAEB) extends embedding evaluation across speech, music, environmental sound, and audio--text tasks, highlighting differences between acoustic and linguistic performance \citep{elassadi2026maeb}. Our focus is complementary: a controlled Sinhala evaluation of speaker identity and transcript-content matching, together with an analysis of how layer selection and task aggregation affect representation recommendations. Rather than assuming that a broad benchmark's overall ranking answers a specific application question, we examine the capabilities underlying that recommendation.

This distinction matters because speech carries several sources of variation simultaneously. Two recordings may be similar because they contain the same words, share a speaker, or originate from similar recording conditions. Labels based on source videos or recording sessions therefore require a different interpretation from verified speaker labels. Likewise, an aggregate containing many speaker tasks can favor speaker-oriented representations even when an application requires content matching. Extraction introduces another choice: the final hidden state is convenient, but its pooled features need not be the most effective representation for every capability. A useful evaluation must distinguish these factors before recommending a checkpoint.

We introduce \textbf{SiMTEB-Audio}, a Sinhala speech-representation evaluation suite comprising fourteen tasks constructed from five corpora. We evaluate twelve frozen encoders spanning wav2vec~2.0/XLS-R, Whisper, WavLM, HuBERT, MMS, and w2v-BERT, alongside a pooled MFCC baseline. The suite includes classification, clustering, pair classification, retrieval, and reranking. We distinguish speaker and transcript-content evaluations from recording-source and quality diagnostics, and audit label provenance, repeated clips, and speaker--prompt relationships. Our content tasks evaluate matching utterances with the same transcript; they do not measure semantic equivalence between differently worded utterances.

The reported results reveal why application-specific evaluation is necessary. Under final-layer extraction, controlled comparisons between a same-transcript utterance from another speaker and a different-transcript utterance from the same speaker expose distinct model preferences. Capability-balanced aggregation also changes the overall ordering. Most importantly, layers selected on OpenSLR-52 development data improve speaker-clustering AMI on OpenSLR-30 for all six encoders in the layer study, with a median absolute gain of 0.572 over their final layers. This cross-corpus result demonstrates that extraction choice can materially change the utility of a frozen checkpoint under the tested protocol. The acoustic baseline is competitive on several default evaluations, reinforcing the need to compare learned representations against simple features rather than attributing every high score to linguistic modeling.

Our contributions are threefold:
\begin{itemize}
    \item A reproducible Sinhala speech-representation evaluation suite with explicit capability definitions, acoustic and chance controls, and separate reporting of recording-related diagnostics.
    \item A controlled analysis of speaker versus transcript-content matching and of layer selection, including development-selected layers evaluated across corpora.
    \item An assessment of recommendation reliability through capability-aware aggregation and cross-corpus comparisons, providing guidance for the evaluated applications rather than a universal best encoder.
\end{itemize}
% Add a MAEB finding only after common task coverage and all intended checkpoints are finalized.
% Verify numerical claims against the final result files and repeated-seed checks before submission.
